Every number in this section is simulated; no number describes a real fleet.

\subsection{Fleet, Parameters and Scenarios}
The fleet has \flShips{} ships in five classes, from an expedition ship of 92
cabins to one of about 2\,750 and a cruise ferry, with the seven references of
Section~\ref{sec:ship}, each made by two makers, one per class. Parameters are
of four kinds; none is an estimate transferred from another source.
\emph{Published}: cabins per class from shipyard fact
sheets~\cite{fincantieri_laperouse,oasis_class}; datasheet predictions at
reference conditions, not observed values, of MTBF 2.19 to 9.09 million hours
for the drivers and supplies (Telcordia SR-332; e.g.~\cite{meanwell_lpv60}) and MTTF
0.54 and 0.56 million hours for the drives (IEC 62380)~\cite{schneider_atv212},
whose inverses serve as nominal \emph{datasheet rates} under a constant-rate
assumption; climate normals; and the mean cruise of 7.0
days~\cite{clia2025}. \emph{Assumed within a published range}: an activation
energy of 0.5\,eV, within MIL-HDBK-217F's values for silicon
microcircuits~\cite{milhdbk217f}. \emph{Assumed, motivated by published
values}: a field factor of 3, below the 5 to 12 by which the handbook raises
naval sheltered over ground benign environments by part class and near a pooled
conversion of 3.3~\cite{jeong2011}. \emph{Assumed}: a bad-lot factor
of 2.5, for which a street-lighting programme whose LED fixtures failed in
their drivers~\cite{halifax2019} shows only that defective driver lots occur;
drivers per cabin
and public area, costs, a Weibull age shape of 1.3, lognormal ship,
ship--reference, maker and cell terms of standard deviation 0.3, 0.3, 0.3 and
0.15, the distribution of cruise lengths (3 to 14 nights, 1 to 2 for the
ferry), and every event and recording parameter. Activation energy and age
shape reach the plans only through temperature and age slopes fitted to the
record. Planned powered hours follow from booked occupancy and climate
normals, failures from realised ones. Of two simulated years, the
first eighteen months of all ships but one per class are fitted and the last
six predicted, at known and at held-out new ships: about \fdMissions{}
ship-missions per fleet.

\emph{Clean} meets the model. \emph{Baseline} adds ship-wide supply events at
0.2 per ship-year, each failing 0.05\,\% of every population; a bad lot
multiplying one reference's rate by 2.5 on half the ships of a class; refits
renewing cabin units; and 10\,\% of preventive replacements recorded as
failures. Other scenarios vary one assumption: truncated Student-$t_3$ latent effects;
25\,\% misrecorded replacements; half the fleet observed for one year;
datasheets that are right (field factor 1); and events at 0.5 per ship-year
failing 0.1\,\%, and at 2 failing 0.05\,\% or 0.2\,\%. The generator draws one
event state per ship-mission, with probability $1-e^{-rH/365}$ for rate $r$
and $H$ nights, and an event adds $\mathrm{Binomial}(N,a)$ failures to every
population; it never places two events in one mission. No public source gives
$r$ or $a$ for events that damage electronics; whole-ship losses of electrical
power on passenger ships, about 0.005 per ship-year from EMSA's 2024 indicator
of 247 occurrences per 1\,000 ship-years times their 2.1\,\% share (our
derivation)~\cite{emsa2025} and 24 in the UK from 2020 to May 2026, without
exposure~\cite{maib_portal}, are context only.

\subsection{Plans and Comparators}
Every plan targets a nominal 99\,\% joint cover for one ship-mission from
pre-sailing information, minimising units under the law used to plan it, and
is scored on realised failures of missions not used to fit or select it; the
term \emph{certified optimum} is reserved for vectors whose search finished.
The \emph{datasheet} benchmark plans for independent Poisson
counts at datasheet rate times planned hours; the \emph{scaled} plan
multiplies those rates by one fleet-wide factor fitted to the record; the
\emph{pooled negative binomial} has covariates and maker but no ship terms.
The \emph{hierarchical} plan uses \brDraws{} predictive draws of the model with
maker, ship and ship--reference terms, and the \emph{fixed-marginal predictive
comparator} the same draws with each reference column permuted independently,
which keeps every empirical marginal and breaks the dependence; it is not the
$f=0$ law of Proposition~\ref{prop:mono} and carries no guarantee.

For draw-based plans the planning law is the empirical law of the draws, and
minimisation is by iterative deepening on the units above the marginal
quantiles. It was certified for all but \fdNotCertJP{}
joint or fixed-marginal and \fdNotCertEv{} incident plans of about
\fdPlansPerType{} per plan type (at least \fdCert\,\% in every group), all
with severe events; these keep the best vector found, feasible under the draws
and at most \fdNotCertGap{} units above a lower bound. With a new seed or
24\,000 draws, only \fdStabSameSeed\,\% and \fdStabSameBig\,\% of vectors on
\fdStabFleets{} fleets were unchanged, yet mean cover stayed at least
\fdStabFresh\,\%: performance is stable where the vector is not.

\begin{table*}[!t]
\caption{Simulated fleet: cover $J$ (\%) and units u per mission of
unit-minimal plans for a nominal 99\,\%; hierarchical (joint) minus fixed-marginal plan from the
same draws, in points ($\Delta J$, 95\,\% interval resampling fleets) and
units ($\Delta$u); italics: $J<99$\,\%, star: interval below 99\,\%
(\textsc{simulation}; \fdFleets{} fleets per scenario)}
\label{tab:fleet}
\centering
\footnotesize
\setlength{\tabcolsep}{1.6pt}
\begin{tabular}{@{}lcccccrlrcrlr@{}}
\toprule
& \multicolumn{8}{c}{known ships} & \multicolumn{4}{c}{new ships} \\
\cmidrule(lr){2-9}\cmidrule(l){10-13}
& datash. & scaled & pooled & hier. & fixed & \multicolumn{2}{c}{$\Delta J$} &
$\Delta$u & hier. & \multicolumn{2}{c}{$\Delta J$} & $\Delta$u \\
Scenario & $J$ & $J$/u & $J$/u & $J$/u & $J$/u & & & & $J$/u & & & \\
\midrule
\tabinput{tables/tab_fleet5.tex}
\bottomrule
\end{tabular}
\end{table*}

\subsection{Rates, Marginals and Dependence}
Table~\ref{tab:fleet} separates three contributions; ranges are over the
scenarios without frequent events. \emph{Learning field rates}: the datasheet
benchmark covers \fdMtbfLo--\fdMtbfHi\,\% of missions at known ships when the
field rate exceeds the datasheet rate, as the design builds in, and
\fdMtbfRight\,\% when datasheets are right; one fleet-wide factor raises this
to \fdScaledLo--\fdScaledHi\,\%. \emph{Estimating each ship's marginal laws}: the pooled,
hierarchical and fixed-marginal plans cover \fdFitKLo--\fdFitKHi\,\% at known
ships and \fdFitNLo--\fdFitNHi\,\% at new ones; in the baseline the
hierarchical plan differs from the pooled negative binomial by \fdJpCov{}
points (\fdJpCovLo{} to \fdJpCovHi) and \fdJpUn{} units (\fdJpUnLo{} to
\fdJpUnHi) per mission at known ships, and by \fdJpNewCov{} points and
\fdJpNewUn{} units at new ones. Joint cover still differs from marginal cover:
per-reference 99\,\% quantiles taken separately cover only
\fdMargLo--\fdMargHi\,\% jointly. \emph{Modelling dependence}: against the
fixed-marginal comparator the hierarchical plan changed cover by \fdDcKLo{} to
\fdDcKHi{} points at known ships and \fdDcNLo{} to \fdDcNHi{} at new ones,
every interval within \fdDcCiKLo{} and \fdDcCiNHi{} points, and saved at most
\fdAbsDuMax{} units per mission, systematically only at new ships. The
permutation left up to \fdResid{} points of residual dependence, the order of
the contrast.

\begin{table}[!t]
\caption{Bridge, large ship: unit-minimal 99\,\% plan with dependence minus
plan under independence, exact and from draws (mean $\pm$ 95\,\% half-width,
\brSeeds{} seeds), scored exactly at the true $f$ (\textsc{exact and
simulation}; $H$
days, $\Lambda$ expected failures, u units at $f=0$)}
\label{tab:bridge}
\centering
\footnotesize
\setlength{\tabcolsep}{2.2pt}
\begin{tabular}{@{}lrrrrrrr@{}}
\toprule
& & & \multicolumn{2}{c}{exact $\Delta$u} & \multicolumn{3}{c}{draws} \\
\cmidrule(lr){4-5}\cmidrule(l){6-8}
$H$ & $\Lambda$ & u & $f{=}0.5$ & $f{=}1$ & $\Delta$u, 0.5 & $\Delta$u, $f{=}1$ &
$\Delta J$, $f{=}1$ \\
\midrule
\tabinput{tables/tab_bridge.tex}
\bottomrule
\end{tabular}
\end{table}

Table~\ref{tab:bridge} links that contrast to the proposition. For the large ship each horizon fixes the marginal
laws, with means proportional to $H$ and a latent multiplier persisting until
resupply, while $f$ varies; plans from \brDraws{} draws of the true law are
set against plans from the same draws permuted. At $f=0$ the draw contrast is Monte Carlo
error, within $\pm$\brNoiseU{} units and $\pm$\brNoiseJ{} points. At the
fleet's spreads, $\sigma=0.42$, and a weekly horizon the exact plans coincide
and the draw contrast is \brFleetDu{} units, as in Table~\ref{tab:fleet}.
Dependence pays when horizon and ship-wide spread grow together: at
$\sigma=0.6$, $f=1$ and 56 days the dependent optimum needs \brBigDep{} units
against \brBigInd, \brBigSave{} or \brBigUnitPct\,\% fewer. These plans
minimise units, as in the fleet; Table~\ref{tab:ship} minimises relative cost.
Plans chosen on draws are nominally 99\,\% and meet it on their draws, but
their mean exact cover was \brJLo--\brJHi\,\%, the error of selecting on a
sample.

\subsection{Supply Events}
Table~\ref{tab:ev} adds to the hierarchical draws an explicit incident term of
the generator's form, one event state per mission adding
$\mathrm{Binomial}(N_p,a)$ failures to every reference, and replans to 99\,\%
joint cover per mission under that law. With the generating $r$ and $a$ it is
an idealised reference; both are also misstated by factors of 0.5 and 2.

\begin{table}[!t]
\caption{Plans with an explicit incident term, known ships: cover (\%), mean
units, units above the plan without the term, units short per event mission
(\textsc{simulation}; \fdFleets{} fleets)}
\label{tab:ev}
\centering
\footnotesize
\setlength{\tabcolsep}{2.4pt}
\begin{tabular}{@{}lcccccc@{}}
\toprule
& \multicolumn{3}{c}{$J$ (\%)} & & & short \\
\cmidrule(lr){2-4}
Plan & all & event & no event & units & extra & /event \\
\midrule
\tabinput{tables/tab_ev.tex}
\bottomrule
\end{tabular}
\end{table}

Without an event term the hierarchical plan covered \fdEvFreqOmitEv\,\%
(\fdEvFreqOmitEvLo--\fdEvFreqOmitEvHi) of event missions at two events a year
failing 0.05\,\%, and \fdEvSevOmitEv\,\% when each fails 0.2\,\%; the
idealised plan raised these to \fdEvFreqKnownEv{} and \fdEvSevKnownEv\,\% and
overall cover from \fdEvFreqOmitAll{} to \fdEvFreqKnownAll{} and from
\fdEvSevOmitAll{} to \fdEvSevKnownAll\,\%, for \fdEvFreqExtra{} and
\fdEvSevExtra\,\% more units. The model of incidents buys this cover and its
parameters size the stock: halving the stated size lost half or more of the
gain on event missions, halving the rate less, and doubling either bought
further cover with stock. In the baseline (\fdEvBaseN{} event missions of
\fdEvBaseAll) it left event-mission cover unchanged (\fdEvBaseOmitEv{} and
\fdEvBaseKnownEv\,\%): events are rarer than the 1\,\% the target allows.
